% GENERATED FILE. Edit canonical lessons, then run npm run generate:books.
% canonical-source-hash: fnv1a64:5466a09795384aad
% canonical-lessons: HI-C534-maic, HI-C534-muqabla, HI-C534-jitana, HI-C534-vijeta, HI-C534-puraskar

\chapter{The Big Match}
\label{ch:hi-the-big-match}
\hlchaptermodality{534}

\begin{chapteropening}
\textbf{By the end of this chapter:} \emph{I can talk about a match and who won it.}

Complete the last lesson of chapter 534: I can talk about a match and who won it.
\end{chapteropening}

\section[\texorpdfstring{\hi{मैच} (maic)}{मैच (maic)}]{\hi{मैच} (maic) — a match (a game)}

\label{lesson:HI-C534-maic}



\begin{quote}
Before the new one: say the Hindi for to pit against, then the Hindi for a stance.
\end{quote}

\subsection*{You'll want to know: \hi{मैच}}
\textbf{\hi{मैच}} — \emph{maic} — \textquotedblleft{}a match (a game)\textquotedblright{}.

\begin{quote}
\hi{कल} \hi{शाम} \hi{को} \hi{क्रिकेट} \hi{का} \hi{मैच} \hi{है।}
\end{quote}

\begin{quote}
\emph{Kal śām ko krikeṭ kā maic hai.}
\end{quote}

\begin{quote}
There is a cricket match tomorrow evening.
\end{quote}

Masculine.

\begin{cousinweb}
\textbf{\hi{मैच}} (\emph{maic}) is a loan from English \emph{match}.
\end{cousinweb}

\begin{grammarlens}[title={What you've built}]
One more word for describing what you lived and felt.
\end{grammarlens}

\subsection*{Your turn}
\begin{itemize}
\raggedright
  \item \emph{Say it:} \emph{maic}
  \item \emph{Say it:} \emph{maic}, once more
  \item \emph{Say it:} \emph{Kal śām ko krikeṭ kā maic hai.} — There is a cricket match tomorrow evening.
  \item \emph{Recall:} say the Hindi for to pit against, then the Hindi for a stance, then say \emph{maic} again
\end{itemize}

\begin{tcolorbox}[breakable,colback=teal!4,colframe=teal!35!black,title={Before you move on}]
What does \hi{मैच} mean? (\textquotedblleft{}A match (a game)\textquotedblright{}.) Say it once more.
\end{tcolorbox}
\section[\texorpdfstring{\hi{मुक़ाबला} (muqāblā)}{मुक़ाबला (muqāblā)}]{\hi{मुक़ाबला} (muqāblā) — a contest; a competition}

\label{lesson:HI-C534-muqabla}



\begin{quote}
Before the new one: say the Hindi for a stance, then the Hindi for a match (a game).
\end{quote}

\subsection*{You'll want to know: \hi{मुक़ाबला}}
\textbf{\hi{मुक़ाबला}} — \emph{muqāblā} — \textquotedblleft{}a contest; a competition\textquotedblright{}.

\begin{quote}
\hi{दोनों} \hi{टीमों} \hi{के} \hi{बीच} \hi{बहुत} \hi{कठिन} \hi{मुक़ाबला} \hi{था।}
\end{quote}

\begin{quote}
\emph{Donõ ṭīmoṁ ke bīc bahut kaṭhin muqāblā thā.}
\end{quote}

\begin{quote}
It was a very hard contest between the two teams.
\end{quote}

Masculine.

\begin{cousinweb}
\textbf{\hi{मुक़ाबला}} (\emph{muqāblā}) is usually traced to Arabic \emph{muqābala}, \textquotedblleft{}facing each other\textquotedblright{}, through Persian.
\end{cousinweb}

\begin{grammarlens}[title={What you've built}]
One more word for describing what you lived and felt.
\end{grammarlens}

\subsection*{Your turn}
\begin{itemize}
\raggedright
  \item \emph{Say it:} \emph{muqāblā}
  \item \emph{Say it:} \emph{muqāblā}, once more
  \item \emph{Say it:} \emph{Donõ ṭīmoṁ ke bīc bahut kaṭhin muqāblā thā.} — It was a very hard contest between the two teams.
  \item \emph{Recall:} say the Hindi for a stance, then the Hindi for a match (a game), then say \emph{muqāblā} again
\end{itemize}

\begin{tcolorbox}[breakable,colback=teal!4,colframe=teal!35!black,title={Before you move on}]
What does \hi{मुक़ाबला} mean? (\textquotedblleft{}A contest; a competition\textquotedblright{}.) Say it once more.
\end{tcolorbox}
\section[\texorpdfstring{\hi{जिताना} (jitānā)}{जिताना (jitānā)}]{\hi{जिताना} (jitānā) — to bring victory to; to win (a match) for}

\label{lesson:HI-C534-jitana}



\begin{quote}
Before the new one: say the Hindi for a match (a game), then the Hindi for a contest.
\end{quote}

\subsection*{You'll want to know: \hi{जिताना}}
\textbf{\hi{जिताना}} — \emph{jitānā} — \textquotedblleft{}to bring victory to; to win (a match) for\textquotedblright{}.

\begin{quote}
\hi{अरुण} \hi{ने} \hi{आख़िर} \hi{में} \hi{दौड़कर} \hi{टीम} \hi{को} \hi{जिताया।}
\end{quote}

\begin{quote}
\emph{Aruṇ ne ākhir meṁ dauṛkar ṭīm ko jitāyā.}
\end{quote}

\begin{quote}
In the end Arun ran it in and won the match for the team.
\end{quote}

Past, with \hi{ने} (he or she …): \hi{उसने} \hi{अपनी} \hi{टीम} \hi{को} \hi{जिताया}, \emph{usne apnī ṭīm ko jitāyā}.

\begin{cousinweb}
\textbf{\hi{जिताना}} (\emph{jitānā}) is the causative of \emph{jītnā}, \textquotedblleft{}to win\textquotedblright{}, which comes from Sanskrit \emph{jita}, \textquotedblleft{}won, conquered\textquotedblright{} (CDIAL 5224): to make someone win.
\end{cousinweb}

\begin{grammarlens}[title={What you've built}]
One more word for describing what you lived and felt.
\end{grammarlens}

\subsection*{Your turn}
\begin{itemize}
\raggedright
  \item \emph{Say it:} \emph{jitānā}
  \item \emph{Say it:} \emph{jitānā}, once more
  \item \emph{Say it:} \emph{Aruṇ ne ākhir meṁ dauṛkar ṭīm ko jitāyā.} — In the end Arun ran it in and won the match for the team.
  \item \emph{Recall:} say the Hindi for a match (a game), then the Hindi for a contest, then say \emph{jitānā} again
\end{itemize}

\begin{tcolorbox}[breakable,colback=teal!4,colframe=teal!35!black,title={Before you move on}]
What does \hi{जिताना} mean? (\textquotedblleft{}To bring victory to; to win (a match) for\textquotedblright{}.) Say it once more.
\end{tcolorbox}
\section[\texorpdfstring{\hi{विजेता} (vijetā)}{विजेता (vijetā)}]{\hi{विजेता} (vijetā) — a winner; a champion}

\label{lesson:HI-C534-vijeta}



\begin{quote}
Before the new one: say the Hindi for a contest, then the Hindi for to bring victory to.
\end{quote}

\subsection*{You'll want to know: \hi{विजेता}}
\textbf{\hi{विजेता}} — \emph{vijetā} — \textquotedblleft{}a winner; a champion\textquotedblright{}.

\begin{quote}
\hi{विजेता} \hi{टीम} \hi{ख़ुशी} \hi{से} \hi{नाच} \hi{रही} \hi{थी।}
\end{quote}

\begin{quote}
\emph{Vijetā ṭīm khuśī se nāc rahī thī.}
\end{quote}

\begin{quote}
The winning team was dancing for joy.
\end{quote}

Masculine or feminine, by the person meant.

\begin{cousinweb}
\textbf{\hi{विजेता}} (\emph{vijetā}) is a tatsama, taken whole from Sanskrit \emph{vijetṛ}, \textquotedblleft{}a vanquisher, a conqueror\textquotedblright{} (Monier-Williams). Hindi takes the Sanskrit nominative form \emph{vijetā}.
\end{cousinweb}

\begin{grammarlens}[title={What you've built}]
One more word for describing what you lived and felt.
\end{grammarlens}

\subsection*{Your turn}
\begin{itemize}
\raggedright
  \item \emph{Say it:} \emph{vijetā}
  \item \emph{Say it:} \emph{vijetā}, once more
  \item \emph{Say it:} \emph{Vijetā ṭīm khuśī se nāc rahī thī.} — The winning team was dancing for joy.
  \item \emph{Recall:} say the Hindi for a contest, then the Hindi for to bring victory to, then say \emph{vijetā} again
\end{itemize}

\begin{tcolorbox}[breakable,colback=teal!4,colframe=teal!35!black,title={Before you move on}]
What does \hi{विजेता} mean? (\textquotedblleft{}A winner; a champion\textquotedblright{}.) Say it once more.
\end{tcolorbox}
\section[\texorpdfstring{\hi{पुरस्कार} (puraskār)}{पुरस्कार (puraskār)}]{\hi{पुरस्कार} (puraskār) — a prize, an award}

\label{lesson:HI-C534-puraskar}



\begin{quote}
Before the new one: say the Hindi for to bring victory to, then the Hindi for a winner.
\end{quote}

\subsection*{You'll want to know: \hi{पुरस्कार}}
\textbf{\hi{पुरस्कार}} — \emph{puraskār} — \textquotedblleft{}a prize, an award\textquotedblright{}.

\begin{quote}
\hi{सबसे} \hi{अच्छे} \hi{खिलाड़ी} \hi{को} \hi{पुरस्कार} \hi{मिलेगा।}
\end{quote}

\begin{quote}
\emph{Sabse acche khilāṛī ko puraskār milegā.}
\end{quote}

\begin{quote}
The best player will get a prize.
\end{quote}

Masculine.

\begin{cousinweb}
\textbf{\hi{पुरस्कार}} (\emph{puraskār}) is a tatsama, taken whole from Sanskrit \emph{puraskāra}, \textquotedblleft{}placing in front, honouring\textquotedblright{} (Monier-Williams). Hindi uses it for a prize or an award.
\end{cousinweb}

\begin{grammarlens}[title={What you've built}]
One more word for describing what you lived and felt.
\end{grammarlens}

\subsection*{Your turn}
\begin{itemize}
\raggedright
  \item \emph{Say it:} \emph{puraskār}
  \item \emph{Say it:} \emph{puraskār}, once more
  \item \emph{Say it:} \emph{Sabse acche khilāṛī ko puraskār milegā.} — The best player will get a prize.
  \item \emph{Recall:} say the Hindi for to bring victory to, then the Hindi for a winner, then say \emph{puraskār} again
\end{itemize}

\begin{tcolorbox}[breakable,colback=teal!4,colframe=teal!35!black,title={Before you move on}]
What does \hi{पुरस्कार} mean? (\textquotedblleft{}A prize, an award\textquotedblright{}.) Say it once more.
\end{tcolorbox}
